% Source: Wald GR, physical PDF p. 18 (printed p. 5).
% Coverage: Figure 1.2, causal structure of spacetime in special relativity.
% Figures/tables/footnotes: Figure only; no tables or footnotes.
% Status: source-reviewed and compile-clean.
% Uncertainties: none.

\begingroup
\renewcommand{\thefigure}{1.2}
\begin{figure}[htbp]
  \centering
  \begin{tikzpicture}[>=Stealth, line cap=round, line join=round]
    \draw[thick] (0,0) -- (-1.80,2.20);
    \draw[thick] (0,0) -- (1.80,2.20);
    \draw[thick] (0,0) -- (-1.80,-2.20);
    \draw[thick] (0,0) -- (1.80,-2.20);
    \draw[thick] (0,2.20) ellipse [x radius=1.80, y radius=0.34];
    \draw[thick] (0,-2.20) ellipse [x radius=1.80, y radius=0.34];
    \fill (0,0) circle (1.7pt) node[left=4pt] {\(p\)};

    \node at (0,3.05) {Future};
    \draw[->, thick, bend right=14] (0.60,2.86) to (0.48,2.30);
    \node at (0,-3.05) {Past};
    \draw[->, thick] (0,-2.78) -- (0,-2.43);

    \node[anchor=west] at (2.35,1.05) {Future light cone};
    \draw[->, thick] (2.20,0.94) -- (1.18,1.38);
    \node[anchor=west] at (2.20,-1.35) {Past light cone};
    \draw[->, thick] (2.05,-1.27) -- (1.12,-1.17);
    \node[anchor=east] at (-2.15,0.10) {Spacelike related};
    % Point into the exterior of the cone, not at its vertex p.
    \draw[->, thick] (-2.00,0.02) -- (-1.05,0.36);
    \draw[->, thick] (-2.00,-0.02) -- (-1.08,-0.38);
  \end{tikzpicture}
  \caption{A diagram showing the causal structure of spacetime in special relativity.
  The ``light cone'' of \(p\) rather than a ``surface of simultaneity'' with \(p\) now
  plays a fundamental role in determining the causal relationship of \(p\) to other
  events.}
  \label{fig:1.2}
\end{figure}
\endgroup
